Deep neural networks (DNNs) are often brittle to small perturbations, which has led to extensive research to certify their robustness. While most existing methods focus on {\em local} robustness, i.e., in the neighborhood of fixed specific inputs, these techniques are often impractical for guaranteeing robustness to {\em real-time} inputs on embedded systems, due to excessive latency and computational cost.

To tackle real-time certification, we proceed in two steps. First, offline, we solve {\em global} robustness problems, which are significantly more complex than local robustness. We then use the results online to perform real-time certification. Offline, we compute {\em bounds} on the difference of output values across different DNN decision classes under both $L_\infty$- {\em and} $L_1$-perturbations. 
{\color{red} On vanila DNNs however, global bounds are deemed to be overly pessimistic, because of inputs that are not In Distribution (ID).
%, in particular 
%Out of Distribution (OOD) inputs. While {\em characterizing} OOD inputs to be removed is impractical, we propose an efficient alternative: 
We propose to project inputs on an AutoEncoder (AE) latent space learnt from ID samples.
With a sufficiently large latent space, ID inputs are only changed slightly and (Linear) AE-DNNs are as accurate as original DNNs, while global bounds are much improved (3-15x).}
To further compute better bounds, we develop novel {\em Diff} MILP models, explicitly representing the small {\em diff}erential variables between the input and its perturbation, instead of implicitly representing them as differences between 2 larger values (standard MILP encoding). 
%To further improve the bounds, we employ principal component analysis (PCA) to focus on {\em realistic} in-distribution inputs. 
Online, our method certifies in real-time over 70\% of incoming images across standard benchmarks (MNIST, Fashion-MNIST, and CIFAR-10 with $L_1$-perturbations of $1, 4, 2$, respectively), requiring only 0.5 ms of latency on a single CPU core.